\documentclass[12pt, a4paper]{article}

% ── Shared preamble (src/lecturenotes.sty) ────────────────────────────────────
%   Packages, colors, boxes, \sol/\soleq and math macros all live in
%   src/*.sty. \input@path points LaTeX two folders up to find them.
\makeatletter
\def\input@path{{../../}}
\makeatother
\usepackage{lecturenotes}

% ── Local TikZ styles (match S1–S4) ───────────────────────────────────────────
\usetikzlibrary{decorations.pathreplacing}
\tikzset{
  layerlabel/.style={font=\small\itshape, gray},
  opbox/.style={draw=orange!70!black, fill=notebg, thick, rounded corners=4pt,
      minimum height=0.9cm, minimum width=1.9cm, font=\small, align=center},
  frame/.style={draw=gray!70, fill=gray!8, thick},
  bbox/.style={line width=1.3pt},
  brace/.style={decorate, decoration={brace, amplitude=4pt}, thick, gray},
}

% \scene{<nut1 style>}{<nut2 style>}{<bolt style>} — toy image: two nuts, one bolt
\newcommand{\scene}[3]{%
  \draw[frame] (0,0) rectangle (3.2,2.6);
  \filldraw[#1] (0.9,1.0) circle (0.45);
  \filldraw[#2] (2.3,1.85) circle (0.4);
  \filldraw[#3] (1.7,0.35) rectangle (2.9,0.75);
}

% \cuboid{<x>}{<depth>}{<height>}{<color>}{<label>} (as in S3/S4)
\newcommand{\cuboid}[5]{%
  \pgfmathsetmacro{\off}{#3*0.3}
  \filldraw[fill=#4!20, draw=#4!80!black] (#1,-#3/2) rectangle ++(#2,#3);
  \filldraw[fill=#4!35, draw=#4!80!black]
  (#1,#3/2) -- ++(\off,\off) -- ++(#2,0) -- ++(-\off,-\off) -- cycle;
  \filldraw[fill=#4!50, draw=#4!80!black]
  (#1+#2,-#3/2) -- ++(\off,\off) -- ++(0,#3) -- ++(-\off,-\off) -- cycle;
  \node[below, font=\scriptsize, align=center] at (#1+#2/2,-1.6) {#5};
}

% ── Document ──────────────────────────────────────────────────────────────────
\begin{document}

\begin{center}
  {\LARGE \textbf{Section 4.5: Beyond Classification}} \\[0.3em]
  {\large\itshape Object Detection, Segmentation and OCR --- an Overview} \\[0.5em]
  {\large AI-IM Class --- Lecture Notes}%
  \ifstudentmode
    \quad {\normalsize\color{gray}[Student Copy]}%
  \fi
\end{center}

\vspace{1em}
\hrule
\vspace{1.5em}

So far a CNN has answered one question per image: \emph{what is in it?}
Many real problems need more. A robot must know \emph{where} each part lies; an
inspection system must outline the scratch, not only report that there is
one; a traceability system must \emph{read} the serial number. The same
convolutional feature extractor (Sections 4.2--4.4) handles all of these. Only
the head and the loss change. This section is an overview; the details of
modern detectors and text recognisers are left to an advanced topic. We
\begin{itemize}[nosep]
  \item compare the main computer-vision tasks,
  \item learn the basic tools of object detection (boxes, IoU, non-maximum
        suppression) and the idea behind YOLO,
  \item see how segmentation predicts a class for every pixel,
  \item outline how OCR reads text, and
  \item map these tasks to manufacturing applications.
\end{itemize}

% ══════════════════════════════════════════════════════════════════════════════
\section{A Family of Vision Tasks}

\begin{figure}[H]
  \centering
  \begin{tikzpicture}[scale=0.95]
    % (a) classification
    \begin{scope}
      \draw[frame] (0,0) rectangle (3.2,2.6);
      \filldraw[fill=gray!45, draw=gray!80!black] (1.6,1.3) circle (0.7);
      \filldraw[fill=gray!8, draw=gray!80!black] (1.6,1.3) circle (0.3);
      \node[font=\small] at (1.6,-0.35) {label: \texttt{nut}};
      \node[font=\small\bfseries] at (1.6,-0.95) {(a) classification};
    \end{scope}
    % (b) detection
    \begin{scope}[xshift=3.7cm]
      \scene{fill=gray!45, draw=gray!80!black}{fill=gray!45, draw=gray!80!black}
      {fill=gray!45, draw=gray!80!black}
      \draw[bbox, mainblue] (0.4,0.5) rectangle (1.4,1.5);
      \draw[bbox, mainblue] (1.85,1.4) rectangle (2.75,2.3);
      \draw[bbox, orange!80!black] (1.6,0.25) rectangle (3.0,0.85);
      \node[font=\tiny, fill=mainblue, text=white, inner sep=1pt, anchor=south west]
      at (0.4,1.5) {nut 0.97};
      \node[font=\tiny, fill=mainblue, text=white, inner sep=1pt, anchor=north east]
      at (1.85,2.3) {nut 0.91};
      \node[font=\tiny, fill=orange!80!black, text=white, inner sep=1pt, anchor=north west]
      at (1.6,0.25) {bolt 0.88};
      \node[font=\small] at (1.6,-0.35) {boxes + labels};
      \node[font=\small\bfseries] at (1.6,-0.95) {(b) object detection};
    \end{scope}
    % (c) semantic segmentation
    \begin{scope}[xshift=7.4cm]
      \scene{fill=mainblue!60, draw=mainblue}{fill=mainblue!60, draw=mainblue}
      {fill=orange!60, draw=orange!80!black}
      \node[font=\small] at (1.6,-0.35) {class of every pixel};
      \node[font=\small\bfseries] at (1.6,-0.95) {(c) semantic segm.};
    \end{scope}
    % (d) instance segmentation
    \begin{scope}[xshift=11.1cm]
      \scene{fill=mainblue!60, draw=mainblue}{fill=green!50!black!50, draw=green!40!black}
      {fill=orange!60, draw=orange!80!black}
      \node[font=\small] at (1.6,-0.35) {pixels of every object};
      \node[font=\small\bfseries] at (1.6,-0.95) {(d) instance segm.};
    \end{scope}
  \end{tikzpicture}
  \caption{Four vision tasks on images of parts. (a) One label for the whole
    image. (b) A box, class and confidence score for every object. (c) Each
    pixel is coloured by class; the two nuts get the same colour. (d) Each
    object gets its own mask, so the two nuts are distinguished.}
  \label{fig:tasks}
\end{figure}

\begin{center}
  \renewcommand{\arraystretch}{1.35}
  \begin{tabular}{lll}
    \toprule
    Task                   & Output per image                                & Loss (main parts)               \\
    \midrule
    Classification         & $C$ class scores                                & cross-entropy                   \\
    Object detection       & a list of (box, class, score)                   & box regression + CE\\
    Semantic segmentation  & $C$ scores per pixel, shape $(C, H, W)$ & CE at every pixel\\
    Instance segmentation  & boxes + a pixel mask per object           & detection + mask loss\\
    \bottomrule
  \end{tabular}
\end{center}

Notice that every loss is built from the two pieces of Section 3.4: a
\emph{regression} loss (MSE or similar) for continuous outputs such as box
coordinates, and \emph{cross-entropy} for class labels.

% ══════════════════════════════════════════════════════════════════════════════
\section{Object Detection}

% ──────────────────────────────────────────────────────────────────────────────
\subsection{Bounding Boxes}

An object's location is a \textbf{bounding box}, the smallest axis-aligned
rectangle around it. Two formats are common (pixel coordinates are measured
from the top-left corner, $y$ pointing down):
\begin{itemize}[nosep]
  \item \textbf{corners} $(x_{\min}, y_{\min}, x_{\max}, y_{\max})$;
  \item \textbf{centre and size} $(x_c, y_c, w, h)$, often divided by the
        image width and height so that all four numbers lie in $[0, 1]$.
        YOLO uses this form.
\end{itemize}

\begin{stepbox}{Exercise 1 \quad Convert a box}
  An image is $200$ pixels wide and $120$ pixels tall. A box has corners
  $(20, 30, 60, 90)$.
  \begin{enumerate}[label=(\alph*), nosep]
    \item In pixels: $x_c = \sol[1cm]{40}$, $y_c = \sol[1cm]{60}$,
          $w = \sol[1cm]{40}$, $h = \sol[1cm]{60}$.
    \item Normalised: $\bigl(\,\sol[1cm]{0.2},\ \sol[1cm]{0.5},\
            \sol[1cm]{0.2},\ \sol[1cm]{0.5}\,\bigr)$.
  \end{enumerate}
\end{stepbox}

% ──────────────────────────────────────────────────────────────────────────────
\subsection{Intersection over Union}

How do we decide whether a predicted box is ``right''? It will never match the
true box exactly. The standard measure of overlap is the
\textbf{intersection over union}:
\begin{resultbox}
  \[
    \text{IoU}(A, B) = \frac{\text{area}(A \cap B)}{\text{area}(A \cup B)}
    = \frac{\text{area}(A \cap B)}
    {\text{area}(A) + \text{area}(B) - \text{area}(A \cap B)} .
  \]
  $\text{IoU} = 1$ for identical boxes and $0$ for boxes that do not touch.
  A detection usually counts as correct if $\text{IoU} \ge 0.5$ with a true
  box of the same class.
\end{resultbox}

\begin{figure}[H]
  \centering
  \begin{tikzpicture}[x=0.7cm, y=0.7cm]
    \draw[gray!30] (0,0) grid (7,6);
    \foreach \x in {0,...,7} \node[below, font=\scriptsize, gray] at (\x,0) {\x};
    \foreach \y in {0,...,6} \node[left, font=\scriptsize, gray] at (0,\y) {\y};
    \fill[purple!30] (2,1) rectangle (4,4);
    \draw[bbox, mainblue] (0,0) rectangle (4,4);
    \draw[bbox, orange!80!black] (2,1) rectangle (6,5);
    \node[font=\small, mainblue] at (1,3.5) {$A$};
    \node[font=\small, orange!80!black] at (5.4,4.5) {$B$};
    \node[font=\scriptsize, purple!70!black] at (3,2.5) {$A\cap B$};
  \end{tikzpicture}
  \caption{Boxes $A = (0, 0, 4, 4)$ and $B = (2, 1, 6, 5)$ in corner format.
    Their intersection is shaded.}
  \label{fig:iou}
\end{figure}

\begin{stepbox}{Exercise 2 \quad Compute IoU}
  \begin{enumerate}[label=(\alph*)]
    \item For the boxes in Figure~\ref{fig:iou}: the intersection is
          $\sol[1cm]{2}$ wide and $\sol[1cm]{3}$ tall, so its area is
          $\sol[1cm]{6}$. The union has area
          $\sol[3cm]{16 + 16 - 6 = 26}$, and
          $\text{IoU} = \sol[2.5cm]{6/26 \approx 0.23}$.
    \item Move $B$ to $(1, 1, 5, 5)$. Now the intersection is
          $\sol[1.2cm]{3 \times 3}$, and
          $\text{IoU} = \sol[3cm]{9/23 \approx 0.39}$.
    \item In general, the intersection of boxes $A$ and $B$ in corner format
          has width
          $\sol[7cm]{\max\bigl(0,\ \min(x^A_{\max}, x^B_{\max}) - \max(x^A_{\min}, x^B_{\min})\bigr)}$,
          and similarly for its height.
  \end{enumerate}
\end{stepbox}

% ──────────────────────────────────────────────────────────────────────────────
\subsection{Non-Maximum Suppression}

A detector looks at many positions, so it usually reports the same object
several times with slightly different boxes. \textbf{Non-maximum suppression
  (NMS)} keeps the best box for each object:

\begin{patternbox}
  \begin{enumerate}[nosep]
    \item Sort all boxes of one class by confidence score.
    \item Keep the box with the highest score.
    \item Delete every remaining box whose IoU with the kept box is above a
          threshold (e.g.\ $0.5$): they are duplicates of it.
    \item Repeat steps 2--3 with the boxes that are left.
  \end{enumerate}
\end{patternbox}

\begin{stepbox}{Exercise 3 \quad Run NMS by hand}
  A detector reports four ``nut'' boxes with scores $A{:}\,0.90$,
  $B{:}\,0.80$, $C{:}\,0.75$, $D{:}\,0.60$. Their pairwise IoU is
  \begin{center}
    \renewcommand{\arraystretch}{1.3}
    \begin{tabular}{c|cccc}
          & $A$  & $B$  & $C$  & $D$  \\
      \hline
      $A$ & ---  & 0.70 & 0.10 & 0.05 \\
      $B$ & 0.70 & ---  & 0.20 & 0.10 \\
      $C$ & 0.10 & 0.20 & ---  & 0.60 \\
      $D$ & 0.05 & 0.10 & 0.60 & ---  \\
    \end{tabular}
  \end{center}
  Use threshold $0.5$.
  \begin{enumerate}[label=(\alph*), nosep]
    \item Round 1: keep $\sol[1cm]{A}$; delete $\sol[2.5cm]{B\ (\text{IoU } 0.70)}$.
    \item Round 2: keep $\sol[1cm]{C}$; delete $\sol[2.5cm]{D\ (\text{IoU } 0.60)}$.
    \item The final detections are $\sol[2cm]{A,\ C}$: the image contains
          $\sol[1cm]{2}$ nuts.
  \end{enumerate}
\end{stepbox}

% ──────────────────────────────────────────────────────────────────────────────
\subsection{Two Families of Detectors}

\begin{itemize}
  \item \textbf{Two-stage detectors} (R-CNN 2014, Fast R-CNN, Faster R-CNN
        2015). Stage 1 proposes a few thousand regions that may contain an
        object; stage 2 runs a CNN classifier and box regressor on each one.
        They are accurate but slower.
  \item \textbf{One-stage detectors} (YOLO 2016, SSD). A single forward pass
        of one network predicts all boxes and classes at once, fast enough for
        live video. Today's YOLO versions are both fast and accurate, which is
        why they are the usual first choice in industry.
\end{itemize}

% ──────────────────────────────────────────────────────────────────────────────
\subsection{The Idea Behind YOLO}

\textbf{YOLO} (``You Only Look Once'', Redmon et al.) turns detection into
a single regression problem. The image is divided into an $S \times S$ grid.
The grid cell that contains the \emph{centre} of an object is responsible
for detecting it. Every cell predicts:
\begin{itemize}[nosep]
  \item $B$ boxes, each with $5$ numbers: $(x_c, y_c, w, h)$ and a
        \emph{confidence} that the box contains an object;
  \item $C$ class probabilities (which class the object is, if there is one).
\end{itemize}

\begin{figure}[H]
  \centering
  \begin{tikzpicture}[x=0.75cm, y=0.75cm]
    \draw[frame] (0,0) rectangle (7,7);
    \fill[mainblue!15] (2,3) rectangle (3,4);
    \draw[gray!60] (0,0) grid (7,7);
    \filldraw[fill=gray!45, draw=gray!80!black] (2.6,3.4) circle (1.1);
    \filldraw[fill=gray!8, draw=gray!80!black] (2.6,3.4) circle (0.45);
    \fill[mainblue!45, opacity=0.55] (2,3) rectangle (3,4);
    \draw[mainblue, thick] (2,3) rectangle (3,4);
    \draw[bbox, mainblue] (1.5,2.3) rectangle (3.7,4.5);
    \fill[red!70!black] (2.6,3.4) circle (2.5pt);
    \node[font=\scriptsize, anchor=west, align=left] at (7.4,3.4)
    {the cell containing the\\object's centre (red dot)\\predicts its box};
    \draw[brace] (0,7.15) -- node[above=4pt, font=\small] {$S = 7$} (7,7.15);
  \end{tikzpicture}
  \caption{YOLO divides the image into an $S \times S$ grid. The shaded cell
    contains the centre of the nut and is responsible for its box.}
  \label{fig:yolo}
\end{figure}

The whole output is therefore one tensor of shape
\begin{resultbox}
  \[
    S \times S \times (5B + C).
  \]
  Afterwards, boxes with low confidence are discarded and NMS removes the
  duplicates.
\end{resultbox}

\begin{stepbox}{Exercise 4 \quad Size of the YOLO output}
  \begin{enumerate}[label=(\alph*), nosep]
    \item The original YOLO used $S = 7$, $B = 2$ and the $C = 20$ classes
          of the PASCAL VOC dataset. Each cell predicts
          $\sol[3cm]{5 \cdot 2 + 20 = 30}$ numbers, and the output has
          $\sol[3cm]{7 \cdot 7 \cdot 30 = 1\,470}$ numbers.
    \item A defect detector uses $S = 13$, $B = 5$ and $C = 3$ defect types.
          The output shape is $\sol[3cm]{13 \times 13 \times 28}$.
  \end{enumerate}
\end{stepbox}

Modern versions refine this idea (several grid sizes for large and small
objects, better box parametrisations, better losses), but the principle is
unchanged: one CNN, one pass, a fixed-size output tensor. In practice
students and engineers use a pretrained model from a library:

\begin{codebox}
from ultralytics import YOLO

model = YOLO("yolov8n.pt")            # small model pretrained on COCO
results = model("conveyor.jpg")      # boxes, classes, scores (after NMS)
model.train(data="defects.yaml", epochs=50)   # fine-tune on our own data
\end{codebox}

% ──────────────────────────────────────────────────────────────────────────────
\subsection{Measuring a Detector}

A predicted box is a \textbf{true positive} if it has $\text{IoU} \ge 0.5$
with a true box of the same class (each true box can be matched only once);
otherwise it is a \textbf{false positive}. True boxes that no prediction
matches are \textbf{missed}.
\[
  \text{precision} = \frac{\text{true positives}}{\text{predicted boxes}},
  \qquad
  \text{recall} = \frac{\text{true positives}}{\text{true boxes}}.
\]
Precision asks ``when the detector says defect, is it right?''; recall asks
``how many of the defects did it find?''.

\begin{stepbox}{Exercise 5 \quad Precision and recall}
  An image set contains $10$ scratches. The detector reports $8$ boxes, of
  which $6$ match a scratch with $\text{IoU} \ge 0.5$.
  \[
    \text{precision} = \sol[2.5cm]{6/8 = 0.75},
    \qquad
    \text{recall} = \sol[2.5cm]{6/10 = 0.60}.
  \]
  Raising the confidence threshold usually makes precision
  $\sol[1.8cm]{\text{higher}}$ and recall $\sol[1.8cm]{\text{lower}}$.
\end{stepbox}

Sweeping the confidence threshold traces a precision--recall curve. The area
under it is the \textbf{average precision} (AP) of a class, and its mean over
all classes is the \textbf{mAP}, the standard single-number score of a
detector.

% ══════════════════════════════════════════════════════════════════════════════
\section{Segmentation}

\textbf{Semantic segmentation} classifies every pixel. For an $H \times W$
image and $C$ classes the network outputs a $(C, H, W)$ tensor of logits:
a softmax over the $C$ classes \emph{at each pixel}, trained with the
cross-entropy averaged over all pixels.

A classification CNN shrinks the feature maps to gain a large receptive
field, but segmentation needs the answer back at full resolution. The
solution is an \textbf{encoder--decoder}: the encoder is an ordinary CNN
that downsamples; the decoder upsamples again (for example with
\emph{transposed convolutions}, which enlarge a feature map). The popular
\textbf{U-Net} (2015) adds skip connections that copy fine detail from each
encoder level to the decoder level of the same size.

\begin{figure}[H]
  \centering
  \begin{tikzpicture}[>=Stealth, scale=0.95, every node/.style={transform shape}]
    \cuboid{0}{0.3}{2.8}{gray}{input\\$(3,H,W)$}
    \cuboid{1.6}{0.5}{2.2}{mainblue}{}
    \cuboid{3.3}{0.8}{1.4}{mainblue}{}
    \cuboid{5.4}{1.1}{0.8}{mainblue}{bottleneck}
    \cuboid{7.8}{0.8}{1.4}{orange}{}
    \cuboid{9.8}{0.5}{2.2}{orange}{}
    \cuboid{11.5}{0.3}{2.8}{green!60!black}{output\\$(C,H,W)$}
    \draw[->, thick] (1.0,0) -- (1.5,0);
    \draw[->, thick] (2.5,0) -- (3.2,0);
    \draw[->, thick] (4.4,0) -- (5.3,0);
    \draw[->, thick] (6.8,0) -- (7.7,0);
    \draw[->, thick] (9.0,0) -- (9.7,0);
    \draw[->, thick] (10.8,0) -- (11.4,0);
    \draw[->, red!70!black, thick, dashed] (2.0,1.7) to[bend left=18] (10.0,1.7);
    \draw[->, red!70!black, thick, dashed] (3.7,1.2) to[bend left=18] (8.2,1.2);
    \node[font=\small, red!70!black] at (6.0,2.9) {skip connections};
    \draw[brace] (4.9,-2.6) -- node[below=5pt, font=\small] {encoder (downsample)} (0,-2.6);
    \draw[brace] (12.4,-2.6) -- node[below=5pt, font=\small] {decoder (upsample)} (7.6,-2.6);
  \end{tikzpicture}
  \caption{An encoder--decoder for segmentation in the style of U-Net. The
    output has one score map per class, at the input resolution.}
  \label{fig:unet}
\end{figure}

\begin{stepbox}{Exercise 6 \quad Segmentation output}
  A crack-inspection network segments $256 \times 256$ images into $C = 3$
  classes (background, crack, stain).
  \begin{enumerate}[label=(\alph*), nosep]
    \item The output tensor for one image has shape
          $\sol[3cm]{(3, 256, 256)}$.
    \item The loss averages the cross-entropy over
          $\sol[3cm]{256^2 = 65\,536}$ pixel predictions.
  \end{enumerate}
\end{stepbox}

\textbf{Instance segmentation} (e.g.\ Mask R-CNN, 2017) combines both ideas:
a detector finds each object, and a small segmentation head predicts a mask
inside each box.

% ══════════════════════════════════════════════════════════════════════════════
\section{Optical Character Recognition (OCR)}

\textbf{OCR} turns text in an image into a string: a serial number on a part,
a lot code and expiry date on a package, the reading of a meter. Modern OCR
systems are a pipeline of the tasks above.

\begin{figure}[H]
  \centering
  \begin{tikzpicture}[>=Stealth, node distance=0.75cm]
    \node[opbox] (img) {image};
    \node[opbox, right=of img] (det) {text\\detection};
    \node[opbox, right=of det] (crop) {crop each\\text line};
    \node[opbox, right=of crop] (rec) {text\\recognition};
    \node[opbox, right=of rec] (post) {post-\\processing};
    \draw[->, thick] (img) -- (det);
    \draw[->, thick] (det) -- (crop);
    \draw[->, thick] (crop) -- (rec);
    \draw[->, thick] (rec) -- (post);
    \node[below=0.35cm of rec, font=\scriptsize\ttfamily] {L0T 2407-B};
    \node[below=0.35cm of post, font=\scriptsize\ttfamily] {LOT 2407-B};
    \node[below=0.35cm of det, font=\scriptsize\itshape, gray, align=center]
    {a detection or\\segmentation CNN};
  \end{tikzpicture}
  \caption{An OCR pipeline. Post-processing applies known rules (here: the
    code starts with the word \texttt{LOT}, so the digit \texttt{0} must be the
    letter \texttt{O}).}
  \label{fig:ocr}
\end{figure}

\begin{enumerate}
  \item \textbf{Text detection} finds the regions that contain text. This is
        object detection or segmentation, as described above, often with
        rotated boxes.
  \item \textbf{Text recognition} reads one cropped line. The classic approach
        cuts the line into single characters and classifies each with a CNN,
        exactly like MNIST digits. This fails when characters touch or are
        unevenly spaced. Modern recognisers such as \textbf{CRNN} instead run
        a CNN over the whole line, read the resulting feature map column by
        column from left to right as a sequence, and predict one character
        (or ``blank'') per column. They are trained with the \textbf{CTC}
        (connectionist temporal classification) loss, which does not need to
        know where each character is.
  \item \textbf{Post-processing} corrects the string with known formats,
        dictionaries or checksums.
\end{enumerate}

\begin{stepbox}{Exercise 7 \quad Decoding a CTC output}
  A CTC recogniser predicts one symbol per column; ``$-$'' is the blank. To
  get the text, first \emph{merge repeated symbols}, then \emph{remove
    blanks}.
  \begin{enumerate}[label=(\alph*), nosep]
    \item \texttt{- 2 2 - 4 4 4 - 0 7 7} \quad $\to$ \quad
          $\sol[3cm]{\texttt{2407}}$
    \item \texttt{- 1 1 - 1 - 5 5} \quad $\to$ \quad
          $\sol[3cm]{\texttt{115}}$
    \item Why is the blank needed? \quad
          $\sol[8cm]{\text{Without it, ``11'' would merge into a single ``1''.}}$
  \end{enumerate}
\end{stepbox}

Ready-made OCR libraries include \textbf{Tesseract}, \textbf{EasyOCR} and
\textbf{PaddleOCR}. For new projects, transformer-based recognisers and large
vision--language models are increasingly used as well.

% ══════════════════════════════════════════════════════════════════════════════
\section{Vision in Manufacturing}

\begin{center}
  \renewcommand{\arraystretch}{1.35}
  \begin{tabular}{>{\raggedright\arraybackslash}p{5.2cm}
      >{\raggedright\arraybackslash}p{3.6cm}
      >{\raggedright\arraybackslash}p{5.4cm}}
    \toprule
    Application                         & Task                    & Example                                         \\
    \midrule
    Accept / reject a part              & classification          & surface defect type on steel sheet              \\
    Locate defects or count parts       & object detection        & missing components on a PCB; parts on a conveyor \\
    Measure defect shape and size       & segmentation            & crack or weld-seam outline                      \\
    Robot picking                       & detection / instance segm. & grasp points in a bin of parts              \\
    Traceability                        & OCR                     & serial numbers, lot codes, expiry dates         \\
    Rare or unknown defects             & anomaly detection       & train only on good parts, flag anything unusual \\
    \bottomrule
  \end{tabular}
\end{center}

\begin{notebox}
  \textbf{Practical advice.} Start from a pretrained model and fine-tune it
  (transfer learning, Section 4.4); a few hundred well-labelled images per
  class are often enough. Invest in consistent lighting and camera position:
  a controlled image makes every task easier. Choose the evaluation measure
  with the cost of errors in mind: a missed defect (low recall) is usually
  far more expensive than a false alarm.
\end{notebox}

% ══════════════════════════════════════════════════════════════════════════════
\section{Summary}

\begin{itemize}
  \item Classification, detection, semantic and instance segmentation share
        a CNN feature extractor; they differ in the head and the loss.
  \item A detection is a box, a class and a score. IoU measures box overlap;
        NMS removes duplicate boxes.
  \item YOLO predicts all boxes in one pass as a tensor of shape
        $S \times S \times (5B + C)$. Two-stage detectors propose regions
        first, then classify them.
  \item Detectors are evaluated by precision, recall and mAP, with matches
        defined by $\text{IoU} \ge 0.5$.
  \item Segmentation predicts a class per pixel with an encoder--decoder such
        as U-Net.
  \item OCR is a pipeline: detect text, recognise each line (CNN features,
        read as a sequence, CTC), then post-process.
\end{itemize}

\paragraph{Advanced topics.}
The YOLO loss function and anchor boxes, the details of CTC, transposed
convolutions and modern segmentation models, and anomaly detection will be
covered in a later topic.

\end{document}
